% 原书分章PDF第10—16页：第26.2—26.3节。
\section{利用高阶TMLE进行推断}

本节假设目标参数在 $P_0$ 处二阶强可微，并且对给定二阶部分典范梯度 $D^{(2)}(P)$，式 (26.4) 在 $P_1=P_0$ 处成立。

\subsection{渐近线性与有效性}

在指定条件下，无论选择哪一个二阶部分典范梯度，2-TMLE都将渐近线性且有效。下述定理精确列出了这些条件。

\begin{bookresult}{定理}{26.1}
假设目标参数 $\Psi$ 具有二阶展开
% 原书行内展开改作无编号展示。
\[
\Psi(P)-\Psi(P_0)=-P_0D^{(1)*}(P)-\frac12P_0^2D^{(2)}(P)+R_3(P,P_0),
\]
且 $P_n^*\in\mathcal M$ 满足方程
\[
P_nD^{(1)*}(P_n^*)=0,\qquad
P_n^2D^{(2)}(P_n^*)=P_n\bar D_n^{(2)}(P_n^*)=0,
\]
其中 $\bar D_n^{(2)}(P_n^*)(o):=n^{-1}\sum_{i=1}^nD^{(2)}(P_n^*)(o,O_i)$。那么，只要
\begin{enumerate}
\item 存在 $P_0$-Donsker类 $\mathcal F$，使 $D^{(1)*}(P_n^*)$ 以趋于一的概率属于 $\mathcal F$，且 $P_0[D^{(1)*}(P_n^*)-D^{(1)*}(P_0)]^2=o_P(1)$；
\item $(P_n^2-P_0^2)D^{(2)}(P_n^*)=o_P(n^{-1/2})$；
\item $R_3(P_n^*,P_0)=o_P(n^{-1/2})$；
\end{enumerate}
$\psi_n^*$ 就是 $\psi_0$ 的渐近线性估计量，其影响函数为 $D^{(1)*}(P_0)$。因此，它也是渐近有效的。
\end{bookresult}

验证这一定理中的条件2，似乎尤其困难。下述引理依赖一致截面变差范数（uniform sectional variation norm；Gill 等，1995）的概念，给出足以确立条件2的一组条件，在实践中可能更容易验证。对给定函数 $f:\mathbb R^d\to\mathbb R$，其一致截面变差范数 $\|f\|_v^*$ 定义为
% 原书行内范数定义改作无编号展示；第二个上确界的下标照录x_s。
\[
\|f\|_v^*:=\sup_s\sup_{x_s}\int|f(dx_s,x_{-s})|,
\]
其中上确界遍及 $f$ 的所有可能截面，而 $\int|f(dx_s,x_{-s})|$ 表示截面 $x_s\mapsto f(x_s,x_{-s})$ 的变差范数。这里，后一截面由给定子集 $s\subseteq\{1,\ldots,d\}$、$x_s:=(x_j:j\in s)$ 及 $x_{-s}:=(x_j:j\notin s)$ 定义。

\begin{bookresult}{引理}{26.1}
如果能够证明：
\begin{enumerate}
\item 映射 $o\mapsto\int D^{(2)}(P_n^*)(o,o_2)\,dP_0(o_2)$ 以及 $o\mapsto\int D^{(2)}(P_n^*)(o_1,o)\,dP_0(o_1)$，以趋于一的概率包含在某个固定的 $P_0$-Donsker类 $\mathcal G$ 中；
\item 下列两个量都依概率趋于零：
\[
\int\left[\int D^{(2)}(P_n^*)(o_1,o_2)\,dP_0(o_1)\right]^2\,dP_0(o_2)
\]
以及
\[
\int\left[\int D^{(2)}(P_n^*)(o_1,o_2)\,dP_0(o_2)\right]^2\,dP_0(o_1);
\]
\item $n^{-1/2}\|D^{(2)}(P_n^*)\|_v^*=o_P(1)$；
\end{enumerate}
那么定理1的条件2成立。
\end{bookresult}

\subsection{置信区间的构造}

基于通常TMLE构造置信区间的技术，同样可用于高阶TMLE。由于2-TMLE $\psi_n^*$ 渐近线性，影响函数为 $D^{(1)*}(P_0)$，因此 $n^{1/2}(\psi_n^*-\psi_0)$ 依分布收敛至均值为零、方差为 $\sigma_0^2:=P_0[D^{(1)*}(P_0)]^2$ 的正态随机变量。这提示Wald型区间
\begin{equation}
\left(\psi_n^*-z_{1-\alpha/2}\sigma_nn^{-1/2},\quad
\psi_n^*+z_{1-\alpha/2}\sigma_nn^{-1/2}\right),
\tag{26.5}
\end{equation}
具有渐近覆盖水平 $(1-\alpha)$。其中，$\sigma_n^2:=P_n[D^{(1)*}(P_n^*)]^2$，$z_\beta$ 为标准正态分布的 $\beta$ 分位数。不同阶的TMLE程序具有相同的一阶行为，但加入高阶项通常能够保证这一行为在更广泛的情形下成立。因此，区间 (26.5) 可以用于 $P_0$ 的任意高阶靶向估计量 $P_n^*$。

上述方法提供渐近正确的推断，但它没有显式利用构造2-TMLE所依据的高阶展开。将这一展开中的高阶项纳入考虑，可能得到有限样本表现更好的置信区间。为此，可以设计一种简单的自助法，其依据是如下事实：如果对 $P_0$ 的某个一致估计量 $P_n^\circ$，有随机样本 $O_1^\#,\ldots,O_n^\#\sim P_n^\circ$，则定义为
\[
\begin{aligned}
&n^{-1/2}\sum_{i=1}^n
\left[D^{(1)*}(P_n^*)(O_i^\#)-P_nD^{(1)*}(P_n^*)\right]\\
&\quad+\frac{n^{-3/2}}2\sum_{i=1}^n\sum_{j=1}^n
\left[D^{(2)}(P_n^*)(O_i^\#,O_j^\#)-P_n^2D^{(2)}(P_n^*)\right]
\end{aligned}
\]
的自助统计量 $Z_n^\#$ 在给定 $P_n$ 时的条件分布，与 $n^{1/2}(\psi_n-\psi_0)$ 的分布，在 $n$ 足够大时，对几乎每个 $P_n$ 都可以彼此任意逼近。$P_n^\circ$ 的显然选择，包括经验测度 $P_n$ 和靶向估计量 $P_n^*$。这提示使用置信区间
\[
\left(\psi_n^*-q_{1-\alpha/2,n}n^{-1/2},\quad
\psi_n^*-q_{\alpha/2,n}n^{-1/2}\right),
\]
其中 $q_{\beta,n}$ 是给定 $P_n$ 时 $Z_n^\#$ 条件分布的 $\beta$ 分位数。这些分位数当然可以通过模拟得到任意精确的估计。这一置信区间易于实施，并且计入了由二阶部分典范梯度 $D^{(2)}(P)$ 所蕴含的V统计量过程解释的二阶变异。

\subsection{高阶TMLE的实现}

高阶TMLE在实践中的实现，并不比普通TMLE更困难，因为前者实际上可以视为后者的一个例子。真正额外的工作在于计算高阶部分典范梯度，或在必要时计算其适当近似，下一节将讨论这一点。下面，我们描述二阶靶向最小损失估计量的实现。

给定来自 $P_0\in\mathcal M$ 的独立随机变量 $O_1,O_2,\ldots,O_n$，我们希望对给定目标参数 $\Psi:\mathcal M\to\mathbb R$，估计 $\psi_0:=\Psi(P_0)$。假定参数 $P\mapsto Q(P)$ 满足：
\begin{enumerate}
\item 对某个映射 $\Psi_1:Q(\mathcal M)\to\mathbb R$，有 $\Psi=\Psi_1\circ Q$；
\item 对某个损失函数 $(o,Q)\mapsto L(Q)(o)$，有 $Q_0=\argmin_{Q\in Q(\mathcal M)}P_0L(Q)$。
\end{enumerate}
这里，$Q_0$ 表示真实汇总量 $Q(P_0)$。为方便记号，下文略微混用符号，以 $\Psi(Q)$ 表示 $\Psi_1(Q)$。可能存在多种参数化方式；通常最好选择其中最简单的一种，使我们能够找到适当损失函数和参数扰动子模型，其在 $Q$ 处的得分张成适当的经验梯度，详见下文。

令 $g=g(P)$ 为干扰参数，记 $g_0:=g(P_0)$。假定 $\Psi$ 的一阶典范梯度 $D^{(1)*}(P)$ 可以表示为 $D^{(1)*}(Q(P),g(P))$，且 $D^{(2)}(P):=D^{(2)}(Q(P),g(P))$ 是 $\Psi$ 的任意相关二阶部分典范梯度。这里也允许类似的记号混用。与前面一样，假定对每个 $o_2$，$o_1\mapsto D^{(2)}(P)(o_1,o_2)$ 属于 $T_{\mathcal M}(P)$。于是，令
\[
\bar D_n^{(2)}(Q,g)(o):=\frac1n\sum_{j=1}^nD^{(2)}(P)(o,O_j),
\]
便有 $\bar D_n^{(2)}(Q,g)\in T_{\mathcal M}(P)$，它将在 $P$ 处发挥二阶得分的作用。

令 $\mathcal Q(Q,g):=\{Q_g(\epsilon):\epsilon\}\subset Q(\mathcal M)$ 为二阶最不利子模型，即 $Q_g(0)=Q$，并且 $D^{(1)*}(Q,g)$ 和 $\bar D_n^{(2)}(Q,g)$ 都属于广义得分向量在 $\epsilon=0$ 处的如下线性张成空间的闭包：
\[
\left\{\left.z^T\frac{\partial}{\partial\epsilon}L(Q_g(\epsilon))\right|_{\epsilon=0}:z\in\mathbb R^p\right\},
\]
其中 $p$ 表示 $\epsilon$ 的维数。

假设我们已有 $(Q_0,g_0)$ 的一个初始估计量 $(Q_{n,0},g_n)$。与普通TMLE一样，2-TMLE通常是迭代程序，尽管在一些重要例子中可以证明只需一步即解析收敛。2-TMLE的一次更新与普通TMLE相同，只是采用二阶而非一阶最不利子模型。具体而言，给定 $Q_0$ 的当前估计 $Q_{n,m}$，更新后的估计 $Q_{n,m+1}$ 定义为
\[
Q_{n,m+1}:=\argmin_{Q\in\mathcal Q(Q_{n,m},g_n)}P_nL(Q).
\]
或者，令 $\epsilon_n^m:=\argmin_\epsilon P_nL(Q_{n,m}(\epsilon))$，便可将 $Q_0$ 的更新估计写成 $Q_{n,m+1}:=Q_{n,m,g_n}(\epsilon_n^m)$。这一迭代更新步骤通常重复至收敛，以 $\epsilon_n^m$ 足够接近零作为判断依据。以 $Q_n^*$ 表示这一迭代程序的极限，$\psi_0$ 的2-TMLE便为 $\psi_n^*:=\Psi(Q_n^*)$。$\psi_n^*$ 的理想渐近性质，在很大程度上来自它按构造满足
\[
P_nD^{(1)*}(Q_n^*,g_n)=P_n\bar D_n^{(2)}(Q_n^*,g_n)=0.
\]

\section{利用近似二阶梯度进行推断}

上面概述的高阶TMLE理论适用于参数具有高阶路径可微性的情形。然而，正如前面指出的，在现实模型中，许多感兴趣的参数作为数据生成分布的泛函，光滑程度不足，甚至不存在二阶梯度。尽管如此，引入近似高阶部分典范梯度后，仍然可以建立有用的高阶理论。

\subsection{渐近线性与有效性}

下述定理是在允许二阶部分典范梯度不存在的情况下，对定理26.1的直接推广；它在2-TMLE的情形下说明了这一点。

\begin{bookresult}{定理}{26.2}
假设对每个 $h$，元素 $D_h^{(2)}(P)\in L_0^2(P^2)$ 满足：对每个 $o\in\mathcal O$，
\[
o_1\mapsto D_h^{(2)}(P)(o_1,o)
\]
属于 $T_{\mathcal M}(P)$；目标参数 $\Psi$ 具有二阶展开
% 原书行内展开改作无编号展示。
\[
\Psi(P)-\Psi(P_0)=-P_0D^{(1)*}(P)-\frac12\lim_{h\to0}P_0^2D_h^{(2)}(P)+R_3(P,P_0),
\]
且 $P_n^*\in\mathcal M$ 满足方程
\[
P_nD^{(1)*}(P_n^*)=0,\qquad
P_n^2D_{h_n}^{(2)}(P_n^*)=P_n\bar D_{h_n,n}^{(2)}(P_n^*)=0,
\]
其中 $\bar D_{h_n,n}^{(2)}(P_n^*)(o):=n^{-1}\sum_{i=1}^nD_{h_n}^{(2)}(P_n^*)(o,O_i)$。那么，只要
\begin{enumerate}
\item 存在 $P_0$-Donsker类 $\mathcal F$，使 $D^{(1)*}(P_n^*)$ 以趋于一的概率属于 $\mathcal F$，且 $P_0[D^{(1)*}(P_n^*)-D^{(1)*}(P_0)]^2=o_P(1)$；
\item $(P_n^2-P_0^2)D_{h_n}^{(2)}(P_n^*)=o_P(n^{-1/2})$；
\item $B_n(h_n):=\frac12[P_0^2D_{h_n}^{(2)}(P_n^*)-\lim_{h\to0}P_0^2D_h^{(2)}(P_n^*)]=o_P(n^{-1/2})$；
\item $R_3(P_n^*,P_0)=o_P(n^{-1/2})$；
\end{enumerate}
$\psi_n^*$ 就是 $\psi_0$ 的渐近线性估计量，其影响函数为 $D^{(1)*}(P_0)$。因此，它也是渐近有效的。
\end{bookresult}

将二阶部分典范梯度替换为其近似后，上述定理条件2的充分条件，与引理26.1中讨论的条件相同。一方面，为适当地控制一致截面变差范数，$h_n$ 不能趋于零过快。另一方面，为使表示误差渐近可忽略，$h_n$ 又必须足够快地趋于零。因此，必须仔细权衡。

\subsection{实现与调节参数的选择}

2-TMLE程序给出 $P_0$ 的一个估计 $P_n^*\in\mathcal M$，使
\[
P_nD^{(1)*}(P_n^*)=P_n\bar D_{h_n,n}^{(2)}(P_n^*)=0,
\]
其中 $\bar D_{h_n,n}^{(2)}(P_n^*)$ 的定义与 $\bar D_n^{(2)}(P_n^*)$ 类似，只是以 $D_{h_n}^{(2)}$ 替换 $D^{(2)}$。因此，可以完全按照第26.2.3节定义2-TMLE程序，只需把来自二阶项的得分替换为相应近似。除了选择适当的调节参数值 $h_n$ 以外，这一算法的实现并不比二阶梯度存在时更困难。关键的实际挑战当然在于确定适当的 $h_n$。

在这一问题中，$h$ 的选择决定了TMLE算法所用的最不利参数子模型。因此，选择 $h$ 完全类似于以下问题中的处理机制或删失机制估计量的选择：在记录了所有潜在混杂因素，或者可用数据受到删失时，估计g计算参数。这些机制是一阶最不利子模型的主要决定因素。van der Laan 和 Gruber（2010）、Gruber 和 van der Laan（2010a）、Stitelman 和 van der Laan（2010）、van der Laan 和 Rose（2011）、Wang 等（2011a）、Gruber 和 van der Laan（2012b）以及本书第10章，对C-TMLE作了描述、详细讨论和实现；它为这一问题提供了有原则的解决方案。同样的方法也可直接用于这里。

C-TMLE是一种在更新步骤中，自动从一组候选最不利参数子模型中进行选择的TMLE算法。对每个候选子模型，所得TMLE算法都会降低经验风险；这一变化的幅度，可以作为评判该子模型价值的标准。C-TMLE恰好利用这一标准，来自适应地构建或选择最不利参数子模型，并由此得到相应TMLE。在某些正则性条件下，只要至少有一个候选子模型允许渐近地完全消除偏差，所得估计量就渐近线性且有效。

具体而言，以 $(Q_{n,0},g_n)$ 表示 $(Q_0,g_0)$ 的初始估计，令 $h$ 索引一个由 $D^{(1)*}(Q_{n,0},g_n)$ 和 $D_h^{(2)}(Q_{n,0},g_n)$ 决定的候选二阶最不利子模型。C-TMLE方案首先选择这样一个 $h$ 值：沿着经过 $Q_{n,0}$ 的相应特定于 $h$ 的参数子模型 $\{Q_{h,n,0}(\epsilon):\epsilon\}$，经验风险 $P_nL(Q_{h,n,0}(\epsilon))$ 可能获得的降幅最大；然后沿这一子模型进行一次更新。换言之，$h_n^0$ 选为映射
\[
h\mapsto P_nL(Q_{h,n,0}(\epsilon_n^0(h)))
\]
的极小点，其中对每个 $h$，$\epsilon_n^0(h):=\argmin_\epsilon P_nL(Q_{h,n,0}(\epsilon))$。第一次更新因而为 $Q_{n,1}:=Q_{h_n^0,n,0}(\epsilon_n^0(h_n^0))$。下一次靶向步骤按类似方式进行，但后续 $h$ 的选择不得大于 $h_n^0$。具体而言，以 $h_n^1$ 表示映射 $h\mapsto P_nL(Q_{h,n,1}(\epsilon_n^1(h)))$ 在区间 $[0,h_n^0]$ 上的极小点，其中对每个 $h$，$\epsilon_n^1(h):=\argmin_\epsilon P_nL(Q_{h,n,1}(\epsilon))$；那么，$Q_0$ 的更新估计为 $Q_{n,2}:=Q_{h_n^1,n,1}(\epsilon_n^1(h_n^1))$。这一迭代过程持续进行，直至按照某个预先指定的标准，例如BIC或某种交叉验证风险，额外一步带来的经验风险下降已不再显著。分别以 $Q_n^*$ 和 $h_n^*$ 表示收敛时 $Q_0$ 的最终估计及 $h_n$ 的最终值。此时，可以直接使用TMLE $\Psi(Q_n^*)$，也可以基于初始估计 $Q_{n,0}$ 和最终调节参数选择 $h=h_n^*$，重新构造2-TMLE。

在上述方案中，特定 $h$ 值的表现是根据一次靶向步骤带来的经验风险降幅来评判的。当TMLE更新过程在一步内即解析收敛时，这似乎是合理的；但在其他情形下，它可能并非最优做法。作为替代，可以对每一个 $h$ 值都将TMLE充分迭代至收敛，在此基础上选择最优 $h$ 值，然后重复整个过程，直至总体收敛。当然，这一算法变体可能需要多得多的计算。

上述方法用于使基于近似二阶部分典范梯度的2-TMLE所需调节参数的选择完全自动化。同样的原则，还可用于构造更复杂的协同TMLE算法；这些算法不仅选择调节参数，还选择决定近似二阶部分典范梯度的其他要素。
